Laws and institutions shape individual outcomes through complex interactions with
citizens' diverse characteristics, yet the mechanisms by which different voting
methods navigate this coupled landscape remain poorly understood. We model
collective governance as optimization on NK fitness landscapes, where a shared
portion of each solution---analogous to laws---is updated through voting, while
an individual portion---analogous to personal traits---remains fixed and unique.
Epistatic interactions between these portions capture how the effects of
legislation depend on individual circumstances. We compare eight voting methods
(plurality, approval, score, Borda, instant-runoff, STAR, minimax, and random
dictator) across a sweep of landscape ruggedness~$K$ and cross-dependency
$\alpha$ between shared and individual portions. We further contrast direct
democracy with representative democracy under varying candidate selection
mechanisms. Our results reveal how voting method choice affects not only
collective fitness but also the distribution of outcomes across the population,
with implications for understanding law as a complex adaptive system.
